% --- CORRECCIÓN IMPORTANTE EN LA PRIMERA LÍNEA ---
% Agregamos "xcolor=table" para que funcione \rowcolor sin errores
\documentclass[aspectratio=169, 10pt, xcolor=table]{beamer}

% --- Configuración del Tema (estilo profesional como el ejemplo) ---
\usetheme{Madrid}
\usecolortheme{dolphin}

% --- Paquetes necesarios ---
\usepackage[utf8]{inputenc}
\usepackage[spanish]{babel}
\usepackage{amsmath, amssymb}
\usepackage{booktabs}
\usepackage{graphicx}
\usepackage{listings}
\usepackage{tikz}
\usetikzlibrary{arrows.meta, positioning, shapes.geometric}

% --- Ruta de Imágenes (crea la carpeta "imagenes" y coloca ahí tus archivos) ---
\graphicspath{{./imagenes/}}

% --- Estilo para código Python ---
\definecolor{codegreen}{rgb}{0,0.6,0}
\definecolor{codegray}{rgb}{0.5,0.5,0.5}
\definecolor{codepurple}{rgb}{0.58,0,0.82}
\definecolor{backcolour}{rgb}{0.95,0.95,0.92}
\lstset{
    language=Python,
    backgroundcolor=\color{backcolour},
    commentstyle=\color{codegreen},
    keywordstyle=\color{blue},
    numberstyle=\tiny\color{codegray},
    stringstyle=\color{codepurple},
    basicstyle=\ttfamily\scriptsize,
    breaklines=true,
    frame=single,
    showstringspaces=false
}

% --- Pie de página personalizado ---
\makeatletter

\setbeamertemplate{footline}{
  \leavevmode%
  \hbox{%
  \begin{beamercolorbox}[wd=.333333\paperwidth,ht=2.25ex,dp=1ex,center]{author in head/foot}%
    \usebeamerfont{author in head/foot}\insertshortauthor
  \end{beamercolorbox}%
  \begin{beamercolorbox}[wd=.333333\paperwidth,ht=2.25ex,dp=1ex,center]{title in head/foot}%
    \usebeamerfont{title in head/foot}Sesión 6
  \end{beamercolorbox}%
  \begin{beamercolorbox}[wd=.333333\paperwidth,ht=2.25ex,dp=1ex,right]{date in head/foot}%
    \usebeamerfont{date in head/foot}Machine Learning \hfill \insertframenumber/\inserttotalframenumber \hspace*{0.3cm}
  \end{beamercolorbox}}%
  \vskip0pt%
}

\makeatother

% --- Datos de la portada ---
\title[SESIÓN 6]{\textbf{Sesión 6}}
\subtitle{Árboles de Decisión y Random Forest}
\author[Prof. C. Sánchez]{Carlos César Sánchez Coronel}
\institute{\textbf{MACHINE LEARNING}}
\date{}

% Logo opcional (descomentar si tienes la imagen)
% \titlegraphic{\includegraphics[height=1.5cm]{logo.png}}

% --- Eliminar la barra de navegación (opcional) ---
\setbeamertemplate{navigation symbols}{}

% --- Índice al inicio de cada sección ---
\AtBeginSection[]{
  \begin{frame}
    \frametitle{Contenido}
    \tableofcontents[currentsection]
  \end{frame}
}

% ============================================================================
% INICIO DEL DOCUMENTO
% ============================================================================
\begin{document}

% --- Portada ---
\begin{frame}
    \titlepage
\end{frame}

% --- Índice general ---
\begin{frame}{Contenido}
    \tableofcontents
\end{frame}

% ============================================================================
% SECCIÓN 1: INTRODUCCIÓN
% ============================================================================
\section{Introducción}

\begin{frame}{Modelos basados en reglas}
    \begin{itemize}
        \item Los árboles de decisión son modelos predictivos que representan reglas en forma jerárquica.
        \item A diferencia de los modelos lineales (que suponen una relación paramétrica global), los árboles \textbf{particionan el espacio} en regiones más simples.
        \item \textbf{Ventaja principal}: interpretabilidad.
        \item \textbf{Problema}: un solo árbol tiene alta varianza (sobreajuste) y es inestable.
        \item \textbf{Solución}: combinar múltiples árboles con \textbf{Random Forest} (ensamble basado en bagging).
    \end{itemize}
    % Basado en introducción de la teoría
    % IMAGEN: Comparación de fronteras lineales vs árbol
\end{frame}

% ============================================================================
% SECCIÓN 2: ÁRBOLES DE DECISIÓN
% ============================================================================
\section{Árboles de Decisión}

\begin{frame}{Estructura de un árbol}
    \begin{columns}[t]
        \column{0.5\textwidth}
        \begin{itemize}
            \item \textbf{Nodo raíz}: contiene todos los datos.
            \item \textbf{Nodos internos}: evalúan una característica y bifurcan.
            \item \textbf{Ramas}: resultados de la evaluación (sí/no, umbral).
            \item \textbf{Nodos hoja}: contienen la predicción final (clase o valor).
            \item \textbf{Profundidad}: número de niveles desde la raíz a una hoja.
        \end{itemize}
        \column{0.5\textwidth}
        % IMAGEN: Diagrama de árbol simple
        \begin{center}
            \begin{tikzpicture}[
                every node/.style={draw, rectangle, rounded corners, align=center, font=\tiny},
                level distance=1cm,
                sibling distance=2.5cm
            ]
                \node {¿Edad $>$ 50?}
                    child { node {¿Ingreso $>$ 50k?}
                        child { node {Clase A} }
                        child { node {Clase B} }
                    }
                    child { node {Clase C} };
            \end{tikzpicture}
        \end{center}
    \end{columns}
    % Basado en Pregunta 1
\end{frame}

\begin{frame}{Criterios de división: Clasificación}
    \begin{columns}[t]
        \column{0.5\textwidth}
        \textbf{Índice de Gini}
        \[
        G = 1 - \sum_{k=1}^{K} p_k^2
        \]
        \begin{itemize}
            \item Mide la probabilidad de clasificar erróneamente un elemento si se etiqueta al azar según la distribución del nodo.
            \item $G=0$: nodo puro.
        \end{itemize}
        \column{0.5\textwidth}
        \textbf{Entropía}
        \[
        H = -\sum_{k=1}^{K} p_k \log_2 p_k
        \]
        \begin{itemize}
            \item Mide el desorden o incertidumbre.
            \item $H=0$: nodo puro; máxima en distribución uniforme.
        \end{itemize}
    \end{columns}
    \vspace{0.3cm}
    \begin{exampleblock}{Ejemplo (60\% A, 40\% B)}
        Gini = $1 - (0.6^2 + 0.4^2) = 0.48$ \\
        Entropía = $-(0.6\log_2 0.6 + 0.4\log_2 0.4) \approx 0.971$
    \end{exampleblock}
    % Basado en Pregunta 2
\end{frame}

\begin{frame}{Ganancia de información}
    \begin{block}{Definición}
        \[
        IG(D, a) = H(D) - \sum_{v \in Valores(a)} \frac{|D_v|}{|D|} H(D_v)
        \]
        Mide la reducción de entropía al dividir según el atributo $a$.
    \end{block}
    \begin{exampleblock}{Ejemplo (Pregunta 3)}
        Nodo padre con $H=0.9$. \\
        División A: hijos con $H=0.3$ (50) y $0.4$ (50) $\Rightarrow IG = 0.9 - (0.5\cdot0.3 + 0.5\cdot0.4) = 0.55$. \\
        División B: hijos con $H=0.1$ (80) y $0.6$ (20) $\Rightarrow IG = 0.9 - (0.8\cdot0.1 + 0.2\cdot0.6) = 0.70$. \\
        Se elige la división B.
    \end{exampleblock}
    % Basado en Pregunta 3
\end{frame}

\begin{frame}{Árboles de regresión}
    \begin{itemize}
        \item Variable objetivo continua.
        \item Criterio de división: minimizar la suma de cuadrados residual (RSS):
        \[
        RSS = \sum_{i \in R_1} (y_i - \bar{y}_{R_1})^2 + \sum_{i \in R_2} (y_i - \bar{y}_{R_2})^2
        \]
        \item Predicción en cada hoja: $\bar{y}$ de las observaciones en esa hoja.
        \item Comparación con regresión lineal: más flexible, pero no extrapola.
    \end{itemize}
    % Basado en Pregunta 4
    % IMAGEN: Árbol de regresión para precios de viviendas
\end{frame}

\begin{frame}{Sobreajuste y poda}
    \begin{alertblock}{Problema}
        Árboles sin restricciones crecen hasta hojas puras → sobreajuste (100\% en entrenamiento, pobre en prueba).
    \end{alertblock}
    \begin{columns}[t]
        \column{0.5\textwidth}
        \textbf{Poda por complejidad-costo}
        \[
        RSS_{\text{podado}} = RSS + \alpha |T|
        \]
        donde $|T|$ = número de hojas, $\alpha$ penaliza la complejidad.
        \column{0.5\textwidth}
        \textbf{Hiperparámetros de parada}
        \begin{itemize}
            \item \texttt{max\_depth}
            \item \texttt{min\_samples\_split}
            \item \texttt{min\_samples\_leaf}
        \end{itemize}
    \end{columns}
    % Basado en Pregunta 5
\end{frame}

\begin{frame}{Ventajas y desventajas de los árboles}
    \begin{table}
        \centering
        \begin{tabular}{p{4.5cm}|p{4.5cm}}
            \toprule
            \textbf{Ventajas} & \textbf{Desventajas} \\
            \midrule
            Interpretables y visualizables & Alta varianza (inestables) \\
            No requieren escalado & Propensos al sobreajuste \\
            Manejan mixtura de variables & Pueden ser sesgados con clases desbalanceadas \\
            Robustos a outliers & En regresión, menos precisos que otros modelos \\
            \bottomrule
        \end{tabular}
    \end{table}
    % Basado en teoría y Pregunta 6
\end{frame}

% ============================================================================
% SECCIÓN 3: RANDOM FOREST
% ============================================================================
\section{Random Forest}

\begin{frame}{Intuición: combinar múltiples árboles}
    \begin{itemize}
        \item Un solo árbol tiene alta varianza.
        \item Si promediamos muchos árboles \textbf{no correlacionados}, la varianza se reduce drásticamente sin aumentar el sesgo.
        \item Random Forest construye una colección de árboles con dos fuentes de aleatoriedad:
        \begin{enumerate}
            \item Muestras bootstrap (bagging).
            \item Selección aleatoria de características en cada división.
        \end{enumerate}
    \end{itemize}
    % Basado en Pregunta 7
    % IMAGEN: Esquema de bagging
\end{frame}

\begin{frame}{Bootstrap aggregating (Bagging)}
    \begin{block}{Procedimiento}
        Para $b = 1,\dots,B$:
        \begin{enumerate}
            \item Generar una muestra bootstrap (con reemplazo) del mismo tamaño que el original.
            \item Entrenar un árbol profundo (sin podar) en esa muestra.
        \end{enumerate}
        Predicción final:
        \begin{itemize}
            \item Clasificación: voto mayoritario.
            \item Regresión: promedio.
        \end{itemize}
    \end{block}
    \begin{itemize}
        \item Aprox. 63.2\% de las observaciones originales aparecen en cada muestra.
        \item El 36.8\% restante se llama \textbf{out-of-bag (OOB)}.
    \end{itemize}
    % Basado en Pregunta 7
\end{frame}

\begin{frame}{Aleatoriedad adicional: selección de características}
    \begin{itemize}
        \item En cada división de cada árbol, solo se considera un subconjunto aleatorio de características (\texttt{max\_features}).
        \item Esto reduce la correlación entre árboles, ya que evita que todos elijan las mismas variables fuertes.
        \item Valores típicos:
        \begin{itemize}
            \item Clasificación: $\sqrt{p}$ (raíz cuadrada del total de características).
            \item Regresión: $p/3$ (un tercio).
        \end{itemize}
    \end{itemize}
    % Basado en Pregunta 8
\end{frame}

\begin{frame}{Error Out-of-Bag (OOB)}
    \begin{block}{Definición}
        Para cada observación $i$, se predice usando solo los árboles donde $i$ no fue incluida en la muestra bootstrap. El error OOB es el error promedio en esas predicciones.
    \end{block}
    \begin{itemize}
        \item Proporciona una estimación del error de prueba \textbf{sin necesidad de validación cruzada}.
        \item Comparable a validación cruzada leave-one-out cuando $B$ es grande.
        \item Precaución: puede ser optimista si hay dependencias (ej. series temporales).
    \end{itemize}
    % Basado en Pregunta 9
\end{frame}

\begin{frame}{Hiperparámetros principales}
    \begin{table}
        \centering
        \begin{tabular}{l|l}
            \toprule
            \textbf{Hiperparámetro} & \textbf{Qué controla} \\
            \midrule
            \texttt{n\_estimators} & Número de árboles (a mayor, más estable) \\
            \texttt{max\_features} & Características por división \\
            \texttt{max\_depth} & Profundidad máxima de cada árbol \\
            \texttt{min\_samples\_split} & Mínimo de muestras para dividir un nodo \\
            \texttt{min\_samples\_leaf} & Mínimo de muestras en una hoja \\
            \texttt{class\_weight} & Peso de clases (para desbalance) \\
            \bottomrule
        \end{tabular}
    \end{table}
    % Basado en Pregunta 10
\end{frame}

\begin{frame}{Importancia de características}
    \begin{columns}[t]
        \column{0.5\textwidth}
        \textbf{Por impureza (MDI)}
        \begin{itemize}
            \item Suma la reducción de impureza (Gini o entropía) aportada por cada variable en todas las divisiones, promediada sobre todos los árboles.
            \item Sesgo hacia variables numéricas o con muchos valores.
        \end{itemize}
        \column{0.5\textwidth}
        \textbf{Por permutación (MDA)}
        \begin{itemize}
            \item Se permuta aleatoriamente la variable en las muestras OOB y se mide el aumento del error.
            \item Más fiable, especialmente con variables correlacionadas.
        \end{itemize}
    \end{columns}
    \vspace{0.3cm}
    \begin{alertblock}{Importancia negativa}
        Si al permutar una variable el error disminuye, significa que esa variable no es importante y posiblemente añade ruido.
    \end{alertblock}
    % Basado en Pregunta 11
\end{frame}

\begin{frame}{Ventajas y desventajas de Random Forest}
    \begin{table}
        \centering
        \begin{tabular}{p{4.5cm}|p{4.5cm}}
            \toprule
            \textbf{Ventajas} & \textbf{Desventajas} \\
            \midrule
            Menor sobreajuste que un solo árbol & Mayor costo computacional \\
            Muy preciso en clasificación/regresión & Pérdida de interpretabilidad global \\
            Maneja alta dimensionalidad & Lento en predicción con muchos árboles \\
            Estima error OOB internamente & No extrapola \\
            Importancia de variables & \\
            Robusto a outliers y faltantes (con imputación) & \\
            \bottomrule
        \end{tabular}
    \end{table}
    % Basado en teoría
\end{frame}

% ============================================================================
% SECCIÓN 4: COMPARACIÓN DE MÉTRICAS
% ============================================================================
\section{Comparación de Métricas}

\begin{frame}{Clasificación: precisión, recall, F1, AUC}
    \begin{itemize}
        \item En problemas desbalanceados (ej. fraude 1\%), la \textbf{exactitud (accuracy)} es engañosa.
        \item Métricas clave:
        \begin{itemize}
            \item \textbf{Recall}: $\frac{VP}{VP+FN}$ (minimizar falsos negativos).
            \item \textbf{Precisión}: $\frac{VP}{VP+FP}$ (minimizar falsos positivos).
            \item \textbf{F1-score}: media armónica de ambas.
            \item \textbf{AUC-ROC}: capacidad discriminativa global.
        \end{itemize}
    \end{itemize}
    % Basado en Pregunta 12
\end{frame}

\begin{frame}{Regresión: RMSE vs MAE}
    \begin{columns}[t]
        \column{0.5\textwidth}
        \textbf{MAE} = $\frac{1}{n}\sum |y_i - \hat{y}_i|$
        \begin{itemize}
            \item Robusto a outliers.
            \item Todos los errores pesan igual.
        \end{itemize}
        \column{0.5\textwidth}
        \textbf{RMSE} = $\sqrt{\frac{1}{n}\sum (y_i - \hat{y}_i)^2}$
        \begin{itemize}
            \item Penaliza más los errores grandes.
            \item Sensible a outliers.
        \end{itemize}
    \end{columns}
    \vspace{0.3cm}
    \begin{exampleblock}{Ejemplo (Pregunta 13)}
        Errores A: [5,5,5,100] → MAE=28.75, RMSE≈50.19 \\
        Errores B: [20,20,20,20] → MAE=20, RMSE=20 \\
        Si los errores grandes son costosos, elegir B.
    \end{exampleblock}
\end{frame}

% ============================================================================
% SECCIÓN 5: CASOS DE USO
% ============================================================================
\section{Casos de Uso}

\begin{frame}{Detección de fraudes (clases desbalanceadas)}
    \begin{itemize}
        \item \textbf{Datos}: transacciones, fraude 0.1\%.
        \item \textbf{Random Forest} con \texttt{class\_weight='balanced'} o SMOTE.
        \item \textbf{Métricas prioritarias}: recall y precisión.
        \item \textbf{Ajuste}: umbral de decisión para equilibrar costos.
    \end{itemize}
    % Basado en Pregunta 12 y 20
\end{frame}

\begin{frame}{Análisis de riesgo crediticio}
    \begin{itemize}
        \item \textbf{Objetivo}: predecir impago.
        \item \textbf{Random Forest} proporciona importancia de variables (historial crediticio, ingresos).
        \item Se pueden visualizar árboles individuales para explicar decisiones a auditores.
    \end{itemize}
    % Basado en Pregunta 11
\end{frame}

\begin{frame}{Predicción de precios de viviendas (regresión)}
    \begin{itemize}
        \item \textbf{Datos}: área, habitaciones, antigüedad, ubicación, etc.
        \item \textbf{Árbol de regresión}: interpretable pero menos preciso.
        \item \textbf{Random Forest Regressor}: mejora precisión, importancia de características.
        \item \textbf{Limitación}: no extrapola a precios muy altos (fuera del rango de entrenamiento).
    \end{itemize}
    % Basado en Pregunta 4 y 17
\end{frame}

% ============================================================================
% SECCIÓN 6: IMPLEMENTACIÓN EN SCIKIT-LEARN
% ============================================================================
\section{Implementación en scikit-learn}

\begin{frame}[fragile]{Árbol de decisión}
    \begin{lstlisting}
from sklearn.tree import DecisionTreeClassifier, plot_tree
import matplotlib.pyplot as plt

clf = DecisionTreeClassifier(max_depth=5, min_samples_split=10,
                             criterion='gini')
clf.fit(X_train, y_train)
y_pred = clf.predict(X_test)

# Visualización
plt.figure(figsize=(20,10))
plot_tree(clf, filled=True, feature_names=feature_names,
          class_names=class_names)
plt.show()

# Importancia de características
importances = clf.feature_importances_
    \end{lstlisting}
\end{frame}

\begin{frame}[fragile]{Random Forest}
    \begin{lstlisting}
from sklearn.ensemble import RandomForestClassifier

rf = RandomForestClassifier(n_estimators=100, max_features='sqrt',
                            oob_score=True, class_weight='balanced',
                            random_state=42)
rf.fit(X_train, y_train)
print(f"OOB Score: {rf.oob_score_}")

# Importancia por impureza
importances = rf.feature_importances_

# Importancia por permutación (más fiable)
from sklearn.inspection import permutation_importance
perm_importance = permutation_importance(rf, X_test, y_test,
                                         n_repeats=10)
    \end{lstlisting}
\end{frame}

% ============================================================================
% SECCIÓN 7: CASO AVANZADO: DETECCIÓN DE FRAUDE (Pregunta 20)
% ============================================================================
\section{Caso Avanzado: Detección de Fraude}

\begin{frame}{Diseño de un sistema antifraude con Random Forest}
    \begin{itemize}
        \item Millones de transacciones/día, fraude 0.1\%.
        \item Requisitos: rapidez en inferencia, interpretabilidad (para analistas).
    \end{itemize}
    \begin{block}{Feature engineering}
        \begin{enumerate}
            \item Log del monto.
            \item Hora del día (codificación cíclica seno/coseno).
            \item Agregadas por usuario: frecuencia últimas 24h, monto promedio, desviación.
            \item Distancia geográfica y tiempo desde última transacción.
            \item Puntuación de anomalía (Isolation Forest) como predictor.
        \end{enumerate}
    \end{block}
\end{frame}

\begin{frame}{Manejo del desbalance y evaluación}
    \begin{itemize}
        \item \textbf{Técnicas}: pesos de clase (\texttt{class\_weight='balanced'}), ajuste de umbral, SMOTE.
        \item \textbf{Métricas}: recall (prioritario), precisión, F1, curva PR.
        \item \textbf{Validación}: usar \texttt{StratifiedKFold} para mantener proporción de clases.
    \end{itemize}
    \vspace{0.3cm}
    \begin{block}{Interpretabilidad}
        \begin{itemize}
            \item Importancia por permutación para explicar variables globales.
            \item SHAP para explicar predicciones individuales a analistas.
        \end{itemize}
    \end{block}
\end{frame}

\begin{frame}{Monitoreo en producción}
    \begin{itemize}
        \item \textbf{Métricas de negocio}: recall $\leftrightarrow$ fraude no detectado (pérdidas), precisión $\leftrightarrow$ transacciones legítimas bloqueadas (insatisfacción).
        \item \textbf{Error OOB} mensual para estimar rendimiento sin validación separada.
        \item \textbf{Detectar drift} en distribuciones de características (PSI, KS).
    \end{itemize}
\end{frame}

% ============================================================================
% SECCIÓN 8: CONCLUSIÓN
% ============================================================================
\section{Conclusión}

\begin{frame}{Lo que hemos aprendido}
    \begin{exampleblock}{Conceptos clave}
        \begin{itemize}
            \item \textbf{Árboles de decisión}: estructura, criterios Gini/entropía, regresión, sobreajuste, poda.
            \item \textbf{Random Forest}: bagging, aleatoriedad en características, error OOB, importancia de variables.
            \item \textbf{Comparación de métricas}: RMSE/MAE en regresión, precisión/recall/F1 en clasificación desbalanceada.
            \item \textbf{Casos de uso}: fraude, riesgo crediticio, predicción de precios.
        \end{itemize}
    \end{exampleblock}

\end{frame}

% --- Diapositiva final ---
\begin{frame}
    \centering
    \Huge \textbf{¡Gracias!}

\end{frame}

\end{document}